% R15 component; only input paths and enumerated location/grammar prose are edited.
% Source: revisions/2026-10-04-r15/manuscript/appendix_observation.tex; commit 1cb3cc9efe135966ec228dfaf51c6841a6dfc97c.
% Generated by code/integrate.py; edit extraction rules there.
\section{Observation and Implementation}\label{app:r15observationtheory}
% Reviewed source: revisions/2026-10-04-r14/manuscript/theory.tex, line 54, commit f5021cefa71492babfcbfa580e0c984f59a9de26
\subsection{Observation and the Economic Implementation}\label{sec:r12observation}
The capital state and the innovation-driven numerical controller are distinct. We now specify an observation model under which the controller uses the history of observed capital rather than latent economic shocks. The model assumes continuous noiseless observation of every capital coordinate, known diffusion and drift coefficients, knowledge of the controller's own actions, and independent private randomization. It does not assume that a finite set of noisy observations reveals Brownian increments.

Write $\Sigma=[\sigma_I I_d,\sigma_C\bm1]$, $Q=\Sigma\Sigma^\top$, and let $n$ be a unit vector spanning the one-dimensional null space of $\Sigma$. Since $\sigma_I>0$, $Q$ is invertible. The observed state path and known actions determine
\[
 M_t=Y_t-Y_0-\int_0^t\{c\bm1+f(Y_s)-a_s\}\,ds.
\]
Let $B^0$ be a scalar Brownian motion independent of the economic shocks and initial state, used only for private randomization.

\begin{proposition}[Implementation from observed state history]\label{prop:r12observation}
Define
\begin{equation}\label{eq:r12recovery}
 \widetilde W_t=\Sigma^\top Q^{-1}M_t+nB^0_t.
\end{equation}
Then $\widetilde W$ is a $(d+1)$-dimensional Brownian motion in the enlarged observation filtration and $\Sigma\widetilde W_t=M_t$. Feeding its cell increments into the prescribed numerical controller, including the componentwise clipping rule, yields the same closed-loop law and payoff as the innovation-driven controller. Thus its existing transfer and policy-specific performance statements apply under this observation model.
\end{proposition}

Private randomization in \eqref{eq:r12recovery} is essential for this exact equivalence: the right-inverse term alone has rank $d$, whereas the implemented clipping acts separately on $d+1$ drivers. The construction does not recover the original unobservable null component path by path. It replaces that component by an independent one with the correct law. Nor is the resulting controller a deterministic function of the state at a single decision date; its memory and independent randomization are part of its definition. The integrated drift in the proposition is an observation-model primitive. A sensor grid, drift quadrature or measurement error would need a further transfer allowance and is not silently treated as exact.


% Reviewed source: revisions/2026-10-04-r14/manuscript/filtration.tex, line 1, commit f5021cefa71492babfcbfa580e0c984f59a9de26
\subsubsection{Information and the comparison class}\label{sec:r14filtration}
The comparison class is defined on a canonical product space. Let
$W=(W^1,\ldots,W^{d+1})$ be the economic Brownian motion and let $R$ be
an independent random element on a standard Borel space. The element $R$
may contain a private Brownian path and a sequence of independent sensor
errors. Write $\mathbb F^W$ for the usual augmentation of the economic
Brownian filtration and $\mathbb F^{W,R}$ for the corresponding product
filtration; the latter reveals only the portion of $R$ specified by its
observation clock. An admissible control is a progressively measurable,
nonanticipative functional of the available paths, taking values in the
primitive action box. Bounded controls and the globally Lipschitz drift
give a unique strong state process and an integrable payoff. These
conventions apply to the comparison class in
Theorems~\ref{thm:r10tube} and~\ref{thm:r11certificate}.

\begin{proposition}[Independent randomization and the optimum]
\label{prop:r14randomization}
Let $V_W^*(y)$ and $V_{W,R}^*(y)$ be the expected-payoff suprema over the
preceding economic and product filtrations. Then
\[
                 V_{W,R}^*(y)=V_W^*(y).
\]
The anchor envelope $G_d$ and the policy-specific regret conclusions
therefore compare the observed-history and finite-observation policies
with the same economic optimum. The equality concerns the supremum;
it neither requires an optimal policy nor asserts that a single fixed
randomizer realization attains that supremum.
\end{proposition}

\begin{proof}
The inequality $V_{W,R}^*\geq V_W^*$ follows by ignoring $R$. For the
reverse inequality, represent an admissible randomized control by its
nonanticipative measurable functional $a(t,W_{[0,t]},R_{[0,t]})$ before
completion of the filtrations. Fixing a full realization $r$ of $R$
leaves a progressively measurable functional of $W$. For almost every
$r$, Fubini's theorem preserves feasibility and the state equation;
strong uniqueness identifies the resulting state with the controlled
state for that section. Independence leaves the law of $W$ equal to
Wiener measure after conditioning on $R=r$. Thus the conditional payoff
is at most $V_W^*(y)$ for almost every $r$. Integrating over $r$ proves
the reverse inequality. Null-set modifications in the augmented
filtration do not change the state law or payoff. Alternatively, the
anchor proof may be repeated directly on the product filtration: its
synchronous comparison and martingale identities use only adaptedness,
bounded actions, and the Brownian property, all of which are preserved.
\end{proof}

\paragraph*{Exact continuous observation.}
For precision in Proposition~\ref{prop:r12observation}, work first in
the product filtration generated by $W$ and an independent scalar
Brownian motion $B^0$. Let $\mathbb G$ be the usual augmentation of
$\sigma(Y_s,B_s^0:0\leq s\leq t)$, with the initial profile included.
The stored numerical state and held actions are functions of this
history. Consequently
\[
 M_t=Y_t-Y_0-\int_0^t\{c\bm1+f(Y_s)-a_s\}\,ds
\]
is $\mathbb G$-adapted. Since it equals $\Sigma W_t$, it is a
$\mathbb G$-martingale: conditional expectation first in the product
filtration and then in its subfiltration $\mathbb G$ proves this claim.
With $A=\Sigma^\top(\Sigma\Sigma^\top)^{-1}$ and a unit null vector
$n$, the process $AM+nB^0$ is a continuous $\mathbb G$-local martingale
starting at zero. Its bracket is $I_{d+1}t$, because
$A\Sigma$ and $nn^\top$ are complementary orthogonal projections and
$[M,B^0]=0$. The multivariate L\'evy characterization gives a Brownian
motion in the stated observation filtration. The identity
$\Sigma(AM+nB^0)=M$ holds pathwise.

This argument also establishes causality. At a decision date, the
already observed drift integral determines the preceding innovation
increment and hence the next internal state. The next action is then
fixed before the next Brownian increment. The held-action state equation
is uniquely solvable on that cell. Induction identifies the joint law
of physical state, internal state, actions, and recovered innovations
with the prescribed innovation controller. The independent null
component supplies the correct law when clipping is applied separately
to $d+1$ coordinates; it does not recover the original latent null
component. The exact observation result is separate from the finite
measurement implementation below.



% Reviewed source: revisions/2026-10-04-r14/manuscript/finite_sensing.tex, line 1, commit f5021cefa71492babfcbfa580e0c984f59a9de26
\subsubsection{Finite observations with a continuous-economy transfer}
\label{sec:r14sensing}
At dates $t_k=kh$, a finite-observation controller receives
$O_k=\overline Y_{t_k}+\eta_k$ and holds
$\overline a_k=\widehat\pi_k(O_k)$ until the next date. The sensor
error is measurable at $t_k$ in the product filtration just specified
and satisfies $\|\eta_k\|_{2,d}\leq\nu_{\rm obs}$, where
$\|X\|_{2,d}=(\E\|X\|_d^2)^{1/2}$. Both the sensor state
$\overline Y$ and the comparison parent's physical state $Y$ obey the
original continuous diffusion. No Brownian coordinate or continuous
drift integral is supplied to the sensor controller.

The protected parent has internal state $Z_0=y$ and selects
$a_k=\widehat\pi_k(Z_k)$. Its update is the specified clipped-innovation
Euler update, with an implementation error of normalized norm at most
$\xi_h$ per cell. This includes innovation conversion and recurrence
arithmetic. Both action maps use the same inward guard and lie within
$\epsilon$ of the common schedule $m_0(t_k)$ coordinatewise.
For a real reference map $\pi_k$, suppose
\[
 \|\pi_k(x)-\pi_k(z)\|_d\leq L\|x-z\|_d,
 \qquad \sup_x\|\widehat\pi_k(x)-\pi_k(x)\|_d\leq\delta_\pi.
\]
Global matrix-norm products and forward error propagation through the
protected activation supply $L$ and $\delta_\pi$ for the saved weights.
Measurement error is measured before the protected map; input
representation within that map is included in $\delta_\pi$.

Retain the production coefficient $\kappa$ of the economic model and
let $\beta\geq\kappa\|B\|_2$. Put
$\overline\sigma=(\sigma_I^2+\sigma_C^2)^{1/2}$ and let $M$ bound
$\|c\bm1+f(y)-m\|_d$ for all feasible held actions. Define the
generator constants $H$ and $K=\beta M+Q_f$ as in the diffusion-transfer argument of Theorem~\ref{thm:r11certificate},
where $H^2\geq d^{-1}\sup_y\|Df(y)\Sigma\|_F^2$, $q_i=(B\Sigma\Sigma^\top B^\top)_{ii}$, and $Q_f=(\kappa/2)(d^{-1}\sum_iq_i^2)^{1/2}$ is the generator
bound, not a diffusion covariance. Let $z>0$ be the clipping threshold and let
$v_{\rm clip}\geq\E(G-\operatorname{clip}(G,-z,z))^2$ for a standard normal
$G$. The bound $2\varphi(z)(z+z^{-1})$ is valid.
Set
\begin{align}
 b_k&=h\{Kt_k/2+H\sqrt{t_k/3}\}
          +\overline\sigma\sqrt{t_kv_{\rm clip}}+k\xi_h,
          \qquad b_0=0,\label{eq:r14sensorforcing}\\
 e_0&=0,\qquad
 a_k^\Delta=\min\{2\epsilon,L(e_k+\nu_{\rm obs})+2\delta_\pi\},\notag\\
 e_{k+1}&=(1+\beta h)e_k+h a_k^\Delta+(b_{k+1}-b_k),
       \qquad 0\leq k<N.\label{eq:r14sensorrecursion}
\end{align}
The forcing is nondecreasing. Each term retains the accumulated
martingale cancellation; in particular, the clipping term is a
finite-horizon second-moment account rather than a sum of cellwise
absolute clipping errors.

Define the nonnegative function
\begin{align*}
 D(t)&=\int_0^t e^{\beta(t-s)}a_{\lfloor s/h\rfloor}^\Delta\,ds,\\
 Q_T&=s_0+(\kappa+\epsilon)T
                    +\sigma_I\sqrt{(1-1/d)T}.
\end{align*}
Here $s_0\geq\|\Pi Y_0\|_{2,d}$, with $\Pi$ the projection away from the common-state direction. For $t\in[t_k,t_{k+1})$, choose $m_{-,k}>0$ below every action in the
common tube and put
\[
 r_k(t)=\epsilon+|m_0(t)-m_0(t_k)|,
 \qquad C_k(t)=r_k(t)
           \{[m_0(t)m_{-,k}]^{-1}+\zeta\}.
\]
With $W(t)=e^{-\rho t}w(t)$ as before, define
\begin{equation}\label{eq:r14sensorpayoff}
 \begin{split}
 B^{\rm sens}_{h,\nu_{\rm obs}}={}&\beta\int_0^T W(t)D(t)\,dt
 +\sum_{k=0}^{N-1}a_k^\Delta
       \int_{t_k}^{t_{k+1}}e^{-\rho t}C_k(t)\,dt\\
 &+2\chi e^{-\rho T}Q_TD(T).
 \end{split}
\end{equation}

\begin{proposition}[Finite-observation payoff transfer]
\label{prop:r14sensor}
Under the stated assumptions,
\begin{align*}
 \|\overline Y_{t_k}-Z_k\|_{2,d}&\leq e_k,\\
 \|\overline a_k-a_k\|_{2,d}&\leq a_k^\Delta,\\
 \|\overline Y_t-Y_t\|_{2,d}&\leq D(t),
\end{align*}
and
\[
       |J_y(\overline a)-J_y(a)|\leq B^{\rm sens}_{h,\nu_{\rm obs}}.
\]
Consequently, a valid improvement interval $[L_a,U_a]$ for the
protected parent at the same time grid transfers to
$[L_a-B^{\rm sens}_{h,\nu_{\rm obs}},U_a+B^{\rm sens}_{h,\nu_{\rm obs}}]$ for the
finite-observation controller. This deterministic operation uses no
additional statistical confidence allocation. It requires a certificate
for that exact protected parent and grid.
\end{proposition}

The allowance concerns the original physical diffusion and uses only
the stated sensor moment bound. A shared-noise fine-Euler comparison
checks the finite-observation implementation and reports its observed
payoff difference. Such a diagnostic does not replace the allowance or
certify the Euler bias. For fixed weights and horizon, uniformly bounded $L,\beta,\epsilon$ and generator constants, and $m_{-,k}\geq\underline m>0$, the forcing vanishes if $h\to0$,
$v_{\rm clip}(h)\to0$, $\xi_h/h\to0$, $\nu_{\rm obs}\to0$, and
$\delta_\pi\to0$. Discrete Gronwall then makes the transfer allowance
vanish. Fixed clipping and fixed numerical or measurement error leave
the corresponding terms in the finite-grid bound.



% Reviewed source: revisions/2026-10-04-r14/manuscript/proofs_new.tex, line 1, commit f5021cefa71492babfcbfa580e0c984f59a9de26
\subsection{Finite-Observation Transfer Proof}\label{app:r14sensingproof}
\subsubsection{Proof of Proposition~\ref{prop:r14sensor}}
All processes are coupled using the same economic Brownian motion and
initial profile; the sensor noise uses the independent randomizer.
The real-map Lipschitz bound, uniform implementation error, and
nonexpansiveness of the common guard give
\[
 \|\widehat\pi_k(\overline Y_{t_k}+\eta_k)
                  -\widehat\pi_k(Z_k)\|_{2,d}
 \leq L\{\|\overline Y_{t_k}-Z_k\|_{2,d}+\nu_{\rm obs}\}+2\delta_\pi.
\]
The common action tube also bounds this difference by $2\epsilon$.

Apply It\^o's formula to $f(\overline Y)$ on each cell, where the
sensor action has already been chosen. The generator and derivative
bounds used in the policy-specific study proof remain valid for this adapted held action.
Stochastic Fubini yields
\begin{align*}
 \int_{t_j}^{t_{j+1}}
       \{f(\overline Y_s)-f(\overline Y_{t_j})\}\,ds
 ={}&\int_{t_j}^{t_{j+1}}(t_{j+1}-v)
       \mathcal A^{\overline a_j}f(\overline Y_v)\,dv\\
 &+\int_{t_j}^{t_{j+1}}(t_{j+1}-v)
             Df(\overline Y_v)\Sigma\,dW_v.
\end{align*}
The $L^2$ norm of the accumulated first term through $t_k$ is at most
$Kht_k/2$. The stochastic integrals are martingale increments over
disjoint cells. It\^o's isometry bounds their accumulated normalized
$L^2$ norm by $Hh\sqrt{t_k/3}$.

Before representation error, every component of
$\Delta W_j-\sqrt h\operatorname{clip}(\Delta W_j/\sqrt h,-z,z)$
has zero mean and variance at most $hv_{\rm clip}$. Components and cells are
independent. Therefore the normalized second moment of the accumulated
projected clipping error is at most
\[
 \frac1d\operatorname{tr}(\Sigma\Sigma^\top)t_kv_{\rm clip}
       =\overline\sigma^2t_kv_{\rm clip}.
\]
The clipping error and the preceding local-drift martingale need not
be independent: the triangle inequality between their already
accumulated norms is sufficient. The implementation error contributes
at most $k\xi_h$.

Let $E_k=\|\overline Y_{t_k}-Z_k\|_{2,d}$. Subtracting the internal
recursion from the integrated sensor state equation thus gives
\[
 E_k\leq\beta h\sum_{j<k}E_j
       +h\sum_{j<k}\|\overline a_j-a_j\|_{2,d}+b_k.
\]
The recursion in~\eqref{eq:r14sensorrecursion} is equivalent to
$e_k=\beta h\sum_{j<k}e_j+h\sum_{j<k}a_j^\Delta+b_k$.
Induction, using the action bound already proved, establishes
$E_k\leq e_k$ and the asserted action inequalities.

The two physical diffusions have the same Brownian term. Subtracting
their integral equations and applying the $\beta$-Lipschitz production
bound gives
\[
 \|\overline Y_t-Y_t\|_{2,d}
 \leq\beta\int_0^t\|\overline Y_s-Y_s\|_{2,d}\,ds
       +\int_0^t a_{\lfloor s/h\rfloor}^\Delta\,ds.
\]
The integral form of Gronwall's inequality proves the bound by $D(t)$.
Projection away from the common-state direction eliminates common
drift, the common schedule, and common diffusion. Bounded production,
the action tube, and It\^o's isometry give, for both physical states,
$\sup_{t\leq T}\|\Pi Y_t\|_{2,d}\leq Q_T$.

It remains to compare payoffs. Subtract the two exact identities
\eqref{eq:r11gain}; the common schedule terms cancel. Production
contributes at most $\beta\int_0^TW(t)D(t)dt$.
For the consumption deficit $R_t$ in~\eqref{eq:r11deficit}, its gradient
in the normalized inner product has components
\[
 -\frac1{m_i}+w(t)+\zeta\bar m
  =\frac1{m_0(t)}-\frac1{m_i}
       +\zeta\{\bar m-m_0(t)\}.
\]
Every point on the segment between the two held actions remains in
their common coordinate tube at $t_k$. Its distance from $m_0(t)$ is
at most $r_k(t)$, and its coordinates are at least $m_{-,k}$. Thus the
normalized gradient norm is bounded by $C_k(t)$. The mean-value
inequality and Cauchy--Schwarz bound the absolute expected deficit
difference by $C_k(t)a_k^\Delta$.

Finally,
\[
 \left|\E\{\|\Pi\overline Y_T\|_d^2-\|\Pi Y_T\|_d^2\}\right|
 \leq\|\overline Y_T-Y_T\|_{2,d}
       \{\|\Pi\overline Y_T\|_{2,d}+\|\Pi Y_T\|_{2,d}\}
 \leq2Q_TD(T).
\]
Adding these three contributions proves
\eqref{eq:r14sensorpayoff}. Adding and subtracting the parent's payoff
on the existing confidence event proves the interval transfer.

\paragraph*{An outward finite computation.}
One need not infer a scalar integral error from a refinement plot.
The function $D$ is nondecreasing and obeys the exact cell recursion
\[
 D(t_{k+1})=e^{\beta h}D(t_k)
                +a_k^\Delta\int_0^h e^{\beta s}\,ds.
\]
An interval enclosure of $m_0$ on each cell gives an upper bound
$C_k^+$ for $C_k(t)$. With the already enclosed weights $A_k,B_k$,
a computable upper allowance is
\[
 \beta\sum_k B_kD(t_{k+1})
       +\sum_k A_kC_k^+a_k^\Delta
       +2\chi e^{-\rho T}Q_TD(T),
\]
evaluated outward. In the forcing recurrence, use nondecreasing saved
upper endpoints $b_k^+\geq b_k$ and the differences of those exact
saved scalars. This preserves the telescoping majorant. The network
and recurrence arithmetic assumptions are the same conditional
binary64 assumptions used for the protected parent.
